% 星际行星（综述）
% license CCBYSA3
% type Wiki

本文根据 CC-BY-SA 协议转载翻译自维基百科\href{https://en.wikipedia.org/wiki/Rogue_planet}{相关文章}。



**流浪行星**（*rogue planet*），也称为**自由漂浮行星**（*free-floating planet*, **FFP**）或**孤立行星质量天体**（*isolated planetary-mass object*, **iPMO**），是指一种具有行星质量、但**不受任何恒星或棕矮星引力束缚**的星际天体。[1][2][3][4]

流浪行星可能起源于其形成所在的行星系统，随后因动力学过程被抛射出系统；它们也可能在**行星系统之外独立形成**。仅银河系中就可能存在**数十亿到数万亿**颗流浪行星，这一数量范围有望由即将发射的**南希·格蕾丝·罗曼空间望远镜**加以精确约束。[5][6] 流浪行星进入太阳系的概率极低，更不用说对地球生命构成直接威胁；在未来 **1000 年**内，其发生概率约为**万亿分之一**。[7]

有些行星质量天体的形成方式可能与恒星相似，国际天文学联合会（IAU）提议将这类天体称为**亚棕矮星**（*sub-brown dwarfs*）。[8] 一个可能的例子是 **Cha 110913−773444**，它要么曾从原行星系统中被抛射而成为流浪行星，要么是独立形成、因此属于亚棕矮星。[9]

---

### 术语

最早的两篇发现论文分别使用了**孤立行星质量天体**（*isolated planetary-mass objects*, iPMO）[10] 和 **自由漂浮行星**（*free-floating planets*, FFP）[11] 这两个名称。大多数天文学论文采用其中之一。[12][13][14]
术语 **“流浪行星”**（*rogue planet*）更常见于**引力微透镜**研究中，而这类研究也经常使用 **FFP** 一词。[15][16]

面向公众的新闻稿则可能采用不同的称呼。例如，2021 年发现至少 **70 颗**自由漂浮行星时，不同新闻稿分别使用了**流浪行星**（rogue planet）[17]、**无恒星行星**（starless planet）[18]、**漂泊行星**（wandering planet）[19] 和 **自由漂浮行星**（free-floating planet）[20] 等名称。
